% ============================================================================
\bigskip
\section{Kết luận và hướng phát triển}
\label{sec:conclusion}
% ============================================================================

\subsection{Kết luận}

Đồ án đã xây dựng hoàn chỉnh một hệ thống đọc hiểu máy trích xuất cho tiếng Việt trên
UIT-ViQuAD~2.0, đáp ứng các yêu cầu của đề tài T11: baseline truy hồi TF-IDF, mô hình
multilingual BERT với lớp QA, XLM-R, đánh giá bằng Exact Match và F1 trên token, và phân tích
theo độ dài đoạn văn cùng phạm vi suy luận của câu hỏi. Gợi ý dùng PhoBERT được thay bằng
ViSoBERT sau khi xác nhận PhoBERT không cung cấp offset mapping cần cho trích xuất; lý do được
trình bày ở Mục~\ref{sec:phobert}. Hệ thống được đóng gói thành ứng dụng web chạy được.

Trên mẫu 500 câu của tập validation, các kết quả chính là:

\begin{enumerate}
  \item \textbf{mBERT fine-tune đạt kết quả tốt nhất: EM 50,80 ($\pm$4,38) và F1 59,49}, vượt
        XLM-R zero-shot (EM 40,60 / F1 56,84) với chênh lệch EM có ý nghĩa thống kê. Baseline
        TF-IDF chỉ đạt EM 0,80 / F1 23,09: so khớp từ khoá tìm được câu chứa đáp án nhưng không
        xác định được đáp án.
  \item \textbf{Điểm tổng che khuất hai kỹ năng tách biệt.} So với XLM-R, mBERT tốt hơn ở cả việc
        xác định chính xác biên đáp án (EM trên câu có đáp án 54,57 so với 45,71) lẫn việc từ
        chối câu không có đáp án (41,01 so với 27,34); trong khi khả năng tìm đúng vùng đáp án, đo
        bằng F1 trên câu có đáp án, tương đương nhau. Fine-tune trong miền chủ yếu dạy mô hình
        \emph{quy ước ranh giới đáp án} và \emph{khi nào nên từ chối}.
  \item \textbf{Lần huấn luyện ViSoBERT đã suy sụp, và đó là một ca lỗi được chẩn đoán chứ không
        phải một phép đo năng lực.} ViSoBERT --- tiền huấn luyện trên văn bản mạng xã hội tiếng
        Việt --- rơi vào hành vi gần như luôn trả rỗng (EM tổng 27,80 bằng đúng tỉ lệ câu không có
        đáp án). Bốn yếu tố đủ để giải thích hiện tượng này mà không cần viện tới chất lượng của
        ViSoBERT như một encoder (Mục~\ref{sec:visobert-factors}): ngân sách \code{max\_length}
        chưa được điều chỉnh cho tokenizer 15.002 token --- \textbf{một nguyên nhân cấu hình chưa
        được bù trừ}, khiến 27,6\% đoạn văn bị cắt thành nhiều cửa sổ so với 2,9\% của mBERT ---
        cùng với hàm mất mát chưa hội tụ, sức chứa nhỏ hơn 45\%, và một cực tiểu địa phương do
        32,39\% câu không có đáp án trong tập huấn luyện tạo ra. Vì vậy \textbf{báo cáo này không
        kết luận rằng tiền huấn luyện tiếng Việt không mang lại lợi ích}; câu hỏi đó vẫn để ngỏ.
  \item \textbf{Câu hỏi ít trùng từ vựng với một câu đơn lẻ khó hơn} với cả bốn mô hình (F1 thấp
        hơn 4,51--11,44 điểm), dù bằng chứng thống kê chỉ ở mức biên. \textbf{Ảnh hưởng của độ dài
        đoạn văn chưa được xác định}: mẫu đánh giá chứa quá ít đoạn văn dài.
\end{enumerate}

Về phương pháp, đồ án cho thấy giá trị thực tế của các cơ chế phòng vệ được cài vào pipeline.
Cổng \code{assert\_gradeable} ngăn việc báo cáo EM vô nghĩa trên test split không có nhãn. Test
cho cửa sổ trượt phát hiện giới hạn 2 cửa sổ âm thầm của thư viện. Hàm lấy mẫu có seed loại bỏ
sai lệch 10 điểm EM do lấy mẫu theo thứ tự tệp. Việc đăng ký giả thuyết trước và các tín hiệu
báo động biến kết quả bất thường của ViSoBERT thành một cuộc điều tra có hệ thống thay vì một con
số bị bỏ qua. Không cơ chế nào trong số này làm con số đẹp hơn; tất cả đều làm con số \emph{đáng
tin} hơn.

\subsection{Giới hạn}
\label{sec:limitations}

Theo yêu cầu của môn học, các giới hạn sau cần được cân nhắc khi đọc mọi kết quả trong báo cáo.

\subsubsection{Giới hạn về dữ liệu và đánh giá}

\begin{enumerate}
  \item \textbf{Đánh giá trên validation, không phải test.} Test split công khai của UIT-ViQuAD~2.0
        không có nhãn (Mục~\ref{sec:blind}). Validation còn được dùng để chọn epoch tốt nhất khi
        fine-tune (trên mẫu 300 câu), nên kết quả của mBERT và ViSoBERT có thể lạc quan nhẹ so với
        một tập kiểm tra độc lập.
  \item \textbf{Mẫu con 500 câu, không phải toàn bộ 3.814 câu validation.} Nửa khoảng tin cậy 95\%
        của EM tổng khoảng $\pm$4 điểm; với các nhóm con 139--218 câu, khoảng này là $\pm$5--8
        điểm. Các chênh lệch nhỏ hơn cỡ đó không được coi là kết luận.
  \item \textbf{Kiểm định thống kê là xấp xỉ.} Các khoảng tin cậy dùng xấp xỉ chuẩn cho tỉ lệ và
        coi hai mô hình như hai mẫu độc lập. Kết quả hiện có chỉ lưu 10 dự đoán mẫu cho mỗi mô hình,
        nên chưa thực hiện được kiểm định cặp (McNemar, paired bootstrap), cũng như khoảng tin cậy
        cho F1.
  \item \textbf{27,8\% câu trong mẫu không có đáp án.} Điểm tổng trộn hai kỹ năng; báo cáo đã tách
        riêng, nhưng mọi so sánh chỉ dựa trên điểm tổng đều cần đọc kèm Bảng~\ref{tab:ans-imp}.
  \item \textbf{Nhãn loại câu hỏi là heuristic tự gán}, không phải nhãn vàng của ViQuAD, và bị
        nhiễu bởi độ trùng từ vựng (Mục~\ref{sec:tagging}). Chỉ 361 câu có đáp án được phân loại.
  \item \textbf{Nhóm độ dài \code{<100} ($n = 1$) và \code{300+} ($n = 9$) không đủ mẫu}; ảnh hưởng
        của độ dài đoạn văn chưa được đo.
  \item \textbf{Chỉ một miền dữ liệu.} UIT-ViQuAD là văn bản Wikipedia với câu hỏi do con người
        viết dựa trên chính đoạn văn. Kết quả không nói lên hiệu năng trên tài liệu học tập, văn
        bản hành chính hay câu hỏi tự nhiên của người dùng thật. Kết quả cao trên benchmark không
        đồng nghĩa với khả năng triển khai thực tế.
\end{enumerate}

\subsubsection{Giới hạn về thiết kế thí nghiệm}

\begin{enumerate}[resume]
  \item \textbf{Mỗi cấu hình chỉ chạy một lần, một seed, không tìm kiếm siêu tham số.} Biến động
        giữa các seed chưa được ước lượng.
  \item \textbf{So sánh giữa các mô hình không phải ablation thuần tuý.} mBERT và ViSoBERT khác nhau
        đồng thời ở ngôn ngữ, miền tiền huấn luyện, từ vựng, tốc độ học, số epoch và $L_{\max}$.
  \item \textbf{$L_{\max} = 30$ cho mBERT và XLM-R} loại khoảng 10,8--12\% đáp án vàng dài hơn 30
        token (đo theo tokenizer của mBERT) khỏi tầm với của mô hình, nên có thể đánh giá thấp nhẹ
        hai mô hình này.
  \item \textbf{PhoBERT chưa được đánh giá}, nên câu hỏi ``encoder tiếng Việt cùng miền có vượt mBERT
        không'' chưa được trả lời; kết luận về ViSoBERT không mở rộng sang PhoBERT.
  \item \textbf{Mức truy vết không đồng đều.} Kết quả chính có tệp JSON kèm commit; các phép đo chẩn
        đoán ViSoBERT (quét ngưỡng, lần chạy đầu) chỉ được ghi trong nhật ký commit. Trường
        \code{commit} ghi commit hiện hành, không ghi nhận thay đổi chưa commit tại thời điểm chạy.
  \item \textbf{Độ trễ phụ thuộc phần cứng} (Apple M5 Pro, MPS) và dao động khoảng 10\% giữa các lần
        chạy. Ngoài ra, số đo trong Bảng~\ref{tab:cost} chạy trên MPS trong khi demo chạy trên CPU:
        đáp án giống hệt nhau, nhưng độ trễ trong container cao hơn.
  \item \textbf{Hai nhóm mô hình được chấm ở hai commit khác nhau} --- \code{590e775} cho mBERT và
        baseline, \code{c63b28c} cho ViSoBERT và XLM-R. Mã đánh giá không thay đổi giữa hai commit
        này, nhưng đây vẫn là một điểm chưa chặt về xuất xứ: một lần chấm lại toàn bộ ở cùng một
        commit sẽ loại bỏ hẳn nghi vấn.
  \item \textbf{Gói demo được sinh từ \code{git archive HEAD}}, nên các tệp chưa commit tại thời
        điểm đóng gói không có trong thư mục \code{source/} của gói. Hai tài liệu thiết kế thuộc
        diện này được đưa kèm riêng trong phụ lục kỹ thuật của bản nộp.
  \item \textbf{ViSoBERT chưa được huấn luyện lại sau khi chẩn đoán.} Lời giải thích ở
        Mục~\ref{sec:visobert-factors} vì vậy chưa được kiểm trực tiếp; nó vẫn là giả thuyết phù
        hợp với bằng chứng, không phải kết luận nhân quả.
\end{enumerate}

\subsection{Hướng phát triển}
\label{sec:future}

Các hướng phát triển được sắp theo mức độ trực tiếp giải quyết các giới hạn ở trên:

\begin{enumerate}
  \item \textbf{Đánh giá toàn bộ validation và lưu dự đoán từng câu.} Chạy \code{run\_eval.py --full}
        và lưu mọi dự đoán để thực hiện kiểm định cặp và tính khoảng tin cậy bootstrap cho cả EM
        lẫn F1. Đây là bước rẻ nhất và nâng độ tin cậy của mọi kết luận.
  \item \textbf{Tách tập dev khỏi train.} Dùng \code{split\_by\_context()} tách một phần train làm tập
        dev để chọn epoch và tinh chỉnh ngưỡng $\tau$, giữ validation hoàn toàn cho đánh giá cuối.
  \item \textbf{Fine-tune XLM-R trên UIT-ViQuAD.} XLM-R zero-shot đã tìm vùng đáp án tốt ngang mBERT
        fine-tune; kết hợp tiền huấn luyện quy mô lớn của XLM-R với dữ liệu trong miền là bước tiếp
        theo tự nhiên và là ứng viên mạnh cho kết quả tốt nhất.
  \item \textbf{Hoàn thiện cấu hình mBERT:} huấn luyện 3--4 epoch (đường cong cho thấy chưa hội tụ),
        nhiều seed, và ablation $L_{\max} \in \{30, 40, 64\}$.
  \item \textbf{Huấn luyện lại ViSoBERT} với \code{max\_length} nâng lên (768) và tốc độ học hạ
        xuống ($3\cdot10^{-5}$), giữ nguyên mọi thứ khác. Đây là phép kiểm trực tiếp và rẻ nhất cho
        lời giải thích ở Mục~\ref{sec:visobert-factors}: nếu ngân sách cửa sổ đúng là yếu tố mạnh
        nhất thì tỉ lệ trả rỗng phải giảm rõ rệt.
  \item \textbf{Hiệu chỉnh ngưỡng $\tau$ trên validation} thay vì để mặc định $0{,}0$. Vì toàn bộ
        lợi thế F1 của mBERT so với XLM-R nằm ở trục từ chối (Bảng~\ref{tab:decomp}), đây là tham
        số rẻ nhất còn lại chưa được tối ưu.
  \item \textbf{Tích hợp PhoBERT} bằng cách tự xây bảng ánh xạ ký tự giữa văn bản đã tách từ và văn bản
        gốc, để trả lời câu hỏi về tiền huấn luyện tiếng Việt khi kiểm soát được miền dữ liệu.
  \item \textbf{Ablation độ dài đoạn văn} với \code{max\_length} $\in \{256, 384, 512\}$ trên một tập
        con được chọn có chủ đích gồm các đoạn văn dài, thay vì dựa vào mẫu ngẫu nhiên.
  \item \textbf{Kiểm chứng heuristic loại câu hỏi} bằng gán nhãn thủ công một mẫu khoảng 200 câu và
        đo độ đồng thuận với heuristic.
  \item \textbf{Phân tích câu hỏi không có đáp án kiểu đối kháng} (thay thực thể, phủ định), nơi cả
        XLM-R và mBERT đều thất bại ở các ví dụ định tính.
  \item \textbf{Hướng triển khai:} lượng tử hoá hoặc xuất ONNX để giảm độ trễ trên CPU; ghép thêm bộ
        truy hồi đoạn văn để hỏi đáp trên kho tài liệu (open-domain QA) --- khi đó cần đánh giá lại
        từ đầu, vì lỗi truy hồi cộng dồn vào lỗi đọc hiểu.
\end{enumerate}
